\section*{Iteración 1}

\begin{table}[H]
\centering
\caption{Caso de prueba de aceptación \#1 Mostrar menú principal (elaboración propia)\label{tab:cpa1}}
\renewcommand{\arraystretch}{1.3}
\begin{tabular}{|p{0.35\textwidth}|p{0.55\textwidth}|}
\hline
\multicolumn{2}{|c|}{\textbf{Caso de Prueba de Aceptación}} \\
\hline
\textbf{Código:} HU\#1\_PA\#1 & \textbf{Historia de Usuario:} 1 \\
\hline
\multicolumn{2}{|l|}{\textbf{Nombre:} Mostrar menú principal} \\
\hline
\multicolumn{2}{|p{0.9\textwidth}|}{\textbf{Descripción:} Prueba realizada para comprobar que el menú principal se muestra correctamente con todos sus botones y componentes.} \\
\hline
\textbf{Pasos de ejecución:} & 1. Se inicia el videojuego. \newline 2. Se verifica que aparezca el menú principal. \newline 3. Se verifica que estén presentes los botones ``Jugar'', ``Instrucciones'' y ``Salir''. \\
\hline
\multicolumn{2}{|c|}{\textbf{Resultado:} Satisfactorio} \\
\hline
\end{tabular}
\end{table}

\begin{table}[H]
\centering
\caption{Caso de prueba de aceptación \#2 Mostrar pantalla de instrucciones (elaboración propia)\label{tab:cpa2}}
\renewcommand{\arraystretch}{1.3}
\begin{tabular}{|p{0.35\textwidth}|p{0.55\textwidth}|}
\hline
\multicolumn{2}{|c|}{\textbf{Caso de Prueba de Aceptación}} \\
\hline
\textbf{Código:} HU\#2\_PA\#1 & \textbf{Historia de Usuario:} 2 \\
\hline
\multicolumn{2}{|l|}{\textbf{Nombre:} Mostrar pantalla de instrucciones} \\
\hline
\multicolumn{2}{|p{0.9\textwidth}|}{\textbf{Descripción:} Prueba realizada para verificar que la pantalla de instrucciones muestra correctamente las reglas del juego.} \\
\hline
\textbf{Pasos de ejecución:} & 1. Se accede al menú principal. \newline 2. Se presiona el botón ``Instrucciones''. \newline 3. Se verifica que se muestre la pantalla con las reglas del juego. \\
\hline
\multicolumn{2}{|c|}{\textbf{Resultado:} Satisfactorio} \\
\hline
\end{tabular}
\end{table}

\begin{table}[H]
\centering
\caption{Caso de prueba de aceptación \#3 Iniciar partida (elaboración propia)\label{tab:cpa3}}
\renewcommand{\arraystretch}{1.3}
\begin{tabular}{|p{0.35\textwidth}|p{0.55\textwidth}|}
\hline
\multicolumn{2}{|c|}{\textbf{Caso de Prueba de Aceptación}} \\
\hline
\textbf{Código:} HU\#3\_PA\#1 & \textbf{Historia de Usuario:} 3 \\
\hline
\multicolumn{2}{|l|}{\textbf{Nombre:} Iniciar partida} \\
\hline
\multicolumn{2}{|p{0.9\textwidth}|}{\textbf{Descripción:} Prueba realizada para verificar que se puede iniciar una nueva partida correctamente.} \\
\hline
\textbf{Pasos de ejecución:} & 1. Se accede al menú principal. \newline 2. Se presiona el botón ``Jugar''. \newline 3. Se verifica que se cargue la pantalla de juego con el tablero visible. \\
\hline
\multicolumn{2}{|c|}{\textbf{Resultado:} Satisfactorio} \\
\hline
\end{tabular}
\end{table}

\begin{table}[H]
\centering
\caption{Caso de prueba de aceptación \#4 Visualizar tablero (elaboración propia)\label{tab:cpa4}}
\renewcommand{\arraystretch}{1.3}
\begin{tabular}{|p{0.35\textwidth}|p{0.55\textwidth}|}
\hline
\multicolumn{2}{|c|}{\textbf{Caso de Prueba de Aceptación}} \\
\hline
\textbf{Código:} HU\#4\_PA\#1 & \textbf{Historia de Usuario:} 4 \\
\hline
\multicolumn{2}{|l|}{\textbf{Nombre:} Visualizar tablero} \\
\hline
\multicolumn{2}{|p{0.9\textwidth}|}{\textbf{Descripción:} Prueba realizada para verificar que el tablero hexagonal se visualiza correctamente con todas las casillas.} \\
\hline
\textbf{Pasos de ejecución:} & 1. Se inicia una partida. \newline 2. Se verifica que el tablero hexagonal se muestre completo. \newline 3. Se verifica que las casillas especiales (escaleras, serpientes, ODS) sean visibles. \\
\hline
\multicolumn{2}{|c|}{\textbf{Resultado:} Satisfactorio} \\
\hline
\end{tabular}
\end{table}

\section*{Iteración 2}

\begin{table}[H]
\centering
\caption{Caso de prueba de aceptación \#5 Lanzar dado (elaboración propia)\label{tab:cpa5}}
\renewcommand{\arraystretch}{1.3}
\begin{tabular}{|p{0.35\textwidth}|p{0.55\textwidth}|}
\hline
\multicolumn{2}{|c|}{\textbf{Caso de Prueba de Aceptación}} \\
\hline
\textbf{Código:} HU\#5\_PA\#1 & \textbf{Historia de Usuario:} 5 \\
\hline
\multicolumn{2}{|l|}{\textbf{Nombre:} Lanzar dado} \\
\hline
\multicolumn{2}{|p{0.9\textwidth}|}{\textbf{Descripción:} Prueba realizada para verificar que el dado se puede lanzar y muestra un resultado aleatorio.} \\
\hline
\textbf{Pasos de ejecución:} & 1. Se inicia una partida. \newline 2. Se presiona el botón de lanzar dado. \newline 3. Se verifica que el dado muestre un número entre 1 y 6. \\
\hline
\multicolumn{2}{|c|}{\textbf{Resultado:} Satisfactorio} \\
\hline
\end{tabular}
\end{table}

\begin{table}[H]
\centering
\caption{Caso de prueba de aceptación \#6 Mover ficha (elaboración propia)\label{tab:cpa6}}
\renewcommand{\arraystretch}{1.3}
\begin{tabular}{|p{0.35\textwidth}|p{0.55\textwidth}|}
\hline
\multicolumn{2}{|c|}{\textbf{Caso de Prueba de Aceptación}} \\
\hline
\textbf{Código:} HU\#6\_PA\#1 & \textbf{Historia de Usuario:} 6 \\
\hline
\multicolumn{2}{|l|}{\textbf{Nombre:} Mover ficha} \\
\hline
\multicolumn{2}{|p{0.9\textwidth}|}{\textbf{Descripción:} Prueba realizada para verificar que la ficha del jugador se mueve correctamente según el resultado del dado.} \\
\hline
\textbf{Pasos de ejecución:} & 1. Se lanza el dado. \newline 2. Se verifica que la ficha se mueva el número de casillas indicado. \newline 3. Se verifica que la ficha llegue a la posición correcta. \\
\hline
\multicolumn{2}{|c|}{\textbf{Resultado:} Satisfactorio} \\
\hline
\end{tabular}
\end{table}

\begin{table}[H]
\centering
\caption{Caso de prueba de aceptación \#7 Aplicar efecto de escalera (elaboración propia)\label{tab:cpa7}}
\renewcommand{\arraystretch}{1.3}
\begin{tabular}{|p{0.35\textwidth}|p{0.55\textwidth}|}
\hline
\multicolumn{2}{|c|}{\textbf{Caso de Prueba de Aceptación}} \\
\hline
\textbf{Código:} HU\#7\_PA\#1 & \textbf{Historia de Usuario:} 7 \\
\hline
\multicolumn{2}{|l|}{\textbf{Nombre:} Aplicar efecto de escalera} \\
\hline
\multicolumn{2}{|p{0.9\textwidth}|}{\textbf{Descripción:} Prueba realizada para verificar que al caer en una casilla con escalera, la ficha avanza automáticamente.} \\
\hline
\textbf{Pasos de ejecución:} & 1. Se mueve la ficha hasta una casilla con escalera. \newline 2. Se verifica que la ficha avance automáticamente a la casilla superior. \newline 3. Se verifica que se muestre la animación correspondiente. \\
\hline
\multicolumn{2}{|c|}{\textbf{Resultado:} Satisfactorio} \\
\hline
\end{tabular}
\end{table}

\begin{table}[H]
\centering
\caption{Caso de prueba de aceptación \#8 Aplicar efecto de serpiente (elaboración propia)\label{tab:cpa8}}
\renewcommand{\arraystretch}{1.3}
\begin{tabular}{|p{0.35\textwidth}|p{0.55\textwidth}|}
\hline
\multicolumn{2}{|c|}{\textbf{Caso de Prueba de Aceptación}} \\
\hline
\textbf{Código:} HU\#8\_PA\#1 & \textbf{Historia de Usuario:} 8 \\
\hline
\multicolumn{2}{|l|}{\textbf{Nombre:} Aplicar efecto de serpiente} \\
\hline
\multicolumn{2}{|p{0.9\textwidth}|}{\textbf{Descripción:} Prueba realizada para verificar que al caer en una casilla con serpiente, la ficha retrocede automáticamente.} \\
\hline
\textbf{Pasos de ejecución:} & 1. Se mueve la ficha hasta una casilla con serpiente. \newline 2. Se verifica que la ficha retroceda automáticamente a la casilla inferior. \newline 3. Se verifica que se muestre la animación correspondiente. \\
\hline
\multicolumn{2}{|c|}{\textbf{Resultado:} Satisfactorio} \\
\hline
\end{tabular}
\end{table}

\section*{Iteración 3}

\begin{table}[H]
\centering
\caption{Caso de prueba de aceptación \#10 Evaluar respuesta (elaboración propia)\label{tab:cpa10}}
\renewcommand{\arraystretch}{1.3}
\begin{tabular}{|p{0.35\textwidth}|p{0.55\textwidth}|}
\hline
\multicolumn{2}{|c|}{\textbf{Caso de Prueba de Aceptación}} \\
\hline
\textbf{Código:} HU\#10\_PA\#1 & \textbf{Historia de Usuario:} 10 \\
\hline
\multicolumn{2}{|l|}{\textbf{Nombre:} Evaluar respuesta} \\
\hline
\multicolumn{2}{|p{0.9\textwidth}|}{\textbf{Descripción:} Prueba realizada para verificar que el sistema evalúa correctamente las respuestas del jugador.} \\
\hline
\textbf{Pasos de ejecución:} & 1. Se responde una pregunta ODS. \newline 2. Se verifica que el sistema indique si la respuesta es correcta o incorrecta. \newline 3. Se verifica que se aplique el efecto correspondiente (avanzar o retroceder). \\
\hline
\multicolumn{2}{|c|}{\textbf{Resultado:} Satisfactorio} \\
\hline
\end{tabular}
\end{table}

\begin{table}[H]
\centering
\caption{Caso de prueba de aceptación \#11 Gestionar turnos (elaboración propia)\label{tab:cpa11}}
\renewcommand{\arraystretch}{1.3}
\begin{tabular}{|p{0.35\textwidth}|p{0.55\textwidth}|}
\hline
\multicolumn{2}{|c|}{\textbf{Caso de Prueba de Aceptación}} \\
\hline
\textbf{Código:} HU\#11\_PA\#1 & \textbf{Historia de Usuario:} 11 \\
\hline
\multicolumn{2}{|l|}{\textbf{Nombre:} Gestionar turnos} \\
\hline
\multicolumn{2}{|p{0.9\textwidth}|}{\textbf{Descripción:} Prueba realizada para verificar que el sistema gestiona correctamente los turnos entre jugadores.} \\
\hline
\textbf{Pasos de ejecución:} & 1. Se completa el turno de un jugador. \newline 2. Se verifica que el turno pase al siguiente jugador. \newline 3. Se verifica que se muestre el indicador del jugador activo. \\
\hline
\multicolumn{2}{|c|}{\textbf{Resultado:} Satisfactorio} \\
\hline
\end{tabular}
\end{table}

\section*{Iteración 4}

\begin{table}[H]
\centering
\caption{Caso de prueba de aceptación \#12 Determinar condición de victoria (elaboración propia)\label{tab:cpa12}}
\renewcommand{\arraystretch}{1.3}
\begin{tabular}{|p{0.35\textwidth}|p{0.55\textwidth}|}
\hline
\multicolumn{2}{|c|}{\textbf{Caso de Prueba de Aceptación}} \\
\hline
\textbf{Código:} HU\#12\_PA\#1 & \textbf{Historia de Usuario:} 12 \\
\hline
\multicolumn{2}{|l|}{\textbf{Nombre:} Determinar condición de victoria} \\
\hline
\multicolumn{2}{|p{0.9\textwidth}|}{\textbf{Descripción:} Prueba realizada para verificar que el sistema detecta correctamente cuando un jugador alcanza la meta.} \\
\hline
\textbf{Pasos de ejecución:} & 1. Se mueve la ficha hasta la casilla final. \newline 2. Se verifica que el sistema detecte la victoria. \newline 3. Se verifica que se muestre la pantalla de victoria. \\
\hline
\multicolumn{2}{|c|}{\textbf{Resultado:} Satisfactorio} \\
\hline
\end{tabular}
\end{table}

\begin{table}[H]
\centering
\caption{Caso de prueba de aceptación \#13 Mostrar estadísticas de partida (elaboración propia)\label{tab:cpa13}}
\renewcommand{\arraystretch}{1.3}
\begin{tabular}{|p{0.35\textwidth}|p{0.55\textwidth}|}
\hline
\multicolumn{2}{|c|}{\textbf{Caso de Prueba de Aceptación}} \\
\hline
\textbf{Código:} HU\#13\_PA\#1 & \textbf{Historia de Usuario:} 13 \\
\hline
\multicolumn{2}{|l|}{\textbf{Nombre:} Mostrar estadísticas de partida} \\
\hline
\multicolumn{2}{|p{0.9\textwidth}|}{\textbf{Descripción:} Prueba realizada para verificar que se muestran correctamente las estadísticas al finalizar la partida.} \\
\hline
\textbf{Pasos de ejecución:} & 1. Se finaliza una partida. \newline 2. Se verifica que se muestren las estadísticas del ganador. \newline 3. Se verifica que se muestren las posiciones de todos los jugadores. \\
\hline
\multicolumn{2}{|c|}{\textbf{Resultado:} Satisfactorio} \\
\hline
\end{tabular}
\end{table}

\begin{table}[H]
\centering
\caption{Caso de prueba de aceptación \#14 Configurar audio (elaboración propia)\label{tab:cpa14}}
\renewcommand{\arraystretch}{1.3}
\begin{tabular}{|p{0.35\textwidth}|p{0.55\textwidth}|}
\hline
\multicolumn{2}{|c|}{\textbf{Caso de Prueba de Aceptación}} \\
\hline
\textbf{Código:} HU\#14\_PA\#1 & \textbf{Historia de Usuario:} 14 \\
\hline
\multicolumn{2}{|l|}{\textbf{Nombre:} Configurar audio} \\
\hline
\multicolumn{2}{|p{0.9\textwidth}|}{\textbf{Descripción:} Prueba realizada para verificar que se puede ajustar el volumen de música y efectos de sonido.} \\
\hline
\textbf{Pasos de ejecución:} & 1. Se accede al menú de configuración. \newline 2. Se ajustan los controles de volumen. \newline 3. Se verifica que los cambios se apliquen correctamente. \\
\hline
\multicolumn{2}{|c|}{\textbf{Resultado:} Satisfactorio} \\
\hline
\end{tabular}
\end{table}

\begin{table}[H]
\centering
\caption{Caso de prueba de aceptación \#15 Pausar partida (elaboración propia)\label{tab:cpa15}}
\renewcommand{\arraystretch}{1.3}
\begin{tabular}{|p{0.35\textwidth}|p{0.55\textwidth}|}
\hline
\multicolumn{2}{|c|}{\textbf{Caso de Prueba de Aceptación}} \\
\hline
\textbf{Código:} HU\#15\_PA\#1 & \textbf{Historia de Usuario:} 15 \\
\hline
\multicolumn{2}{|l|}{\textbf{Nombre:} Pausar partida} \\
\hline
\multicolumn{2}{|p{0.9\textwidth}|}{\textbf{Descripción:} Prueba realizada para verificar que se puede pausar y reanudar la partida.} \\
\hline
\textbf{Pasos de ejecución:} & 1. Durante una partida, se presiona el botón de pausa. \newline 2. Se verifica que el juego se detenga. \newline 3. Se verifica que se pueda reanudar la partida. \\
\hline
\multicolumn{2}{|c|}{\textbf{Resultado:} Satisfactorio} \\
\hline
\end{tabular}
\end{table}

\begin{table}[H]
\centering
\caption{Caso de prueba de aceptación \#16 Reiniciar partida (elaboración propia)\label{tab:cpa16}}
\renewcommand{\arraystretch}{1.3}
\begin{tabular}{|p{0.35\textwidth}|p{0.55\textwidth}|}
\hline
\multicolumn{2}{|c|}{\textbf{Caso de Prueba de Aceptación}} \\
\hline
\textbf{Código:} HU\#16\_PA\#1 & \textbf{Historia de Usuario:} 16 \\
\hline
\multicolumn{2}{|l|}{\textbf{Nombre:} Reiniciar partida} \\
\hline
\multicolumn{2}{|p{0.9\textwidth}|}{\textbf{Descripción:} Prueba realizada para verificar que se puede reiniciar una partida en curso.} \\
\hline
\textbf{Pasos de ejecución:} & 1. Durante una partida, se accede al menú de pausa. \newline 2. Se selecciona la opción ``Reiniciar''. \newline 3. Se verifica que la partida se reinicie con todas las fichas en la posición inicial. \\
\hline
\multicolumn{2}{|c|}{\textbf{Resultado:} Satisfactorio} \\
\hline
\end{tabular}
\end{table}

\section*{Iteración 5}

\begin{table}[H]
\centering
\caption{Caso de prueba de aceptación \#17 Habilitar cámara de seguimiento (elaboración propia)\label{tab:cpa17}}
\renewcommand{\arraystretch}{1.3}
\begin{tabular}{|p{0.35\textwidth}|p{0.55\textwidth}|}
\hline
\multicolumn{2}{|c|}{\textbf{Caso de Prueba de Aceptación}} \\
\hline
\textbf{Código:} HU\#17\_PA\#1 & \textbf{Historia de Usuario:} 17 \\
\hline
\multicolumn{2}{|l|}{\textbf{Nombre:} Habilitar cámara de seguimiento} \\
\hline
\multicolumn{2}{|p{0.9\textwidth}|}{\textbf{Descripción:} Prueba realizada para verificar que la cámara sigue automáticamente a la ficha del jugador activo.} \\
\hline
\textbf{Pasos de ejecución:} & 1. Se mueve una ficha en el tablero. \newline 2. Se verifica que la cámara se desplace siguiendo la ficha. \newline 3. Se verifica que la ficha permanezca visible en todo momento. \\
\hline
\multicolumn{2}{|c|}{\textbf{Resultado:} Satisfactorio} \\
\hline
\end{tabular}
\end{table}

\begin{table}[H]
\centering
\caption{Caso de prueba de aceptación \#18 Habilitar modo de práctica (elaboración propia)\label{tab:cpa18}}
\renewcommand{\arraystretch}{1.3}
\begin{tabular}{|p{0.35\textwidth}|p{0.55\textwidth}|}
\hline
\multicolumn{2}{|c|}{\textbf{Caso de Prueba de Aceptación}} \\
\hline
\textbf{Código:} HU\#18\_PA\#1 & \textbf{Historia de Usuario:} 18 \\
\hline
\multicolumn{2}{|l|}{\textbf{Nombre:} Habilitar modo de práctica} \\
\hline
\multicolumn{2}{|p{0.9\textwidth}|}{\textbf{Descripción:} Prueba realizada para verificar que el modo de práctica permite responder preguntas sin afectar el progreso en el tablero.} \\
\hline
\textbf{Pasos de ejecución:} & 1. Se accede al modo de práctica desde el menú principal. \newline 2. Se responden preguntas ODS. \newline 3. Se verifica que no se apliquen efectos de avance o retroceso. \\
\hline
\multicolumn{2}{|c|}{\textbf{Resultado:} Satisfactorio} \\
\hline
\end{tabular}
\end{table}
